% =====================================================================
\section{Звук: пищалка и мелодии}
% =====================================================================
\lessonintro
  {извлекать звук заданной высоты, понимать связь частоты и ноты и проигрывать мелодию, не останавливая программу}
  {урок 4 (ШИМ и частота), урок 21 (отсчёт времени без \code{delay()})}
  {плата подаёт звуковой сигнал по кнопке и играет мелодию, оставаясь готовой к другим делам}
  {пассивная пьезопищалка (пьезоизлучатель)}

\subsection{Сначала: звук --- это быстрое мигание}

Светодиод, мигающий 440 раз в секунду, глаз видит как ровный свет. Пищалка, на которую подали такой же
сигнал, 440 раз в секунду толкает воздух --- и ухо слышит ноту «ля». Чем чаще колебания, тем выше звук.
Значит, для звука не нужно ничего нового: достаточно менять уровень на выводе с нужной частотой.

\subsection{Частота и нота}

\begin{figure}[H]
\centering
\begin{tikzpicture}[x=1cm,y=0.8cm,font=\sffamily\scriptsize]
  \node[left] at (0,2.9) {262 Гц («до»)};
  \draw[very thick,accent] (0,2.4) \foreach \k in {0,...,3}{ -- (\k*2.4,3.4) -- (\k*2.4+1.2,3.4) -- (\k*2.4+1.2,2.4) -- (\k*2.4+2.4,2.4)};
  \node[left] at (0,0.9) {523 Гц («до» октавой выше)};
  \draw[very thick,esp] (0,0.4) \foreach \k in {0,...,7}{ -- (\k*1.2,1.4) -- (\k*1.2+0.6,1.4) -- (\k*1.2+0.6,0.4) -- (\k*1.2+1.2,0.4)};
  \draw[<->] (0,3.75) -- node[above]{период $T = 1/262 \approx 3{,}8$ мс} (2.4,3.75);
  \draw[->] (0,-0.1) -- (10.2,-0.1) node[right]{$t$};
\end{tikzpicture}
\caption{Прямоугольный сигнал двух частот. Удвоение частоты повышает звук ровно на октаву}
\end{figure}

Соседние ноты отличаются по частоте в одно и то же число раз --- примерно в 1,06. Двенадцать таких шагов
дают удвоение частоты, то есть октаву. За точку отсчёта принята нота «ля» первой октавы: \SI{440}{Hz}.

\begin{center}\small
\begin{tabular}{@{}lcccccccc@{}}
\toprule
Нота & до & ре & ми & фа & соль & ля & си & до \\
Обозначение & C4 & D4 & E4 & F4 & G4 & A4 & B4 & C5 \\
\midrule
Частота, Гц & 262 & 294 & 330 & 349 & 392 & 440 & 494 & 523 \\
\bottomrule
\end{tabular}
\end{center}

\subsection{Какая бывает пищалка}

\begin{center}\small
\begin{tabularx}{\textwidth}{@{}l X X@{}}
\toprule
 & Пассивная (пьезоизлучатель) & Активная (со встроенным генератором) \\
\midrule
Что подавать & сигнал нужной частоты & просто HIGH \\
Высота звука & любая, задаёт программа & одна, задана на заводе \\
Как отличить & при подаче 3,3 В только щёлкает & при подаче 3,3 В пищит \\
\bottomrule
\end{tabularx}
\end{center}

Для мелодий нужна \textbf{пассивная} пищалка. Пьезоизлучатель потребляет очень малый ток, поэтому его
подключают прямо между \pin{21} и GND. Электромагнитную пищалку и тем более динамик так подключать
нельзя: им нужен транзисторный ключ из урока 16.

Для опытов урока понадобится ещё кнопка из урока 3: она включена между \pin{5} и GND, подтяжку
к питанию включает программа.

\begin{figure}[H]
\centering
\begin{circuitikz}[scale=0.95,transform shape]
  \node[blk,fill=blkcpu,minimum width=2.8cm,minimum height=4.2cm] (m) at (0,0) {};
  \node[font=\sffamily\small\bfseries] at (0,2.6) {ESP32-S3};
  \foreach \y/\n in {1/{\pin{21}},-0.5/{\pin{5}},-1.7/GND}
    \node[font=\sffamily\scriptsize,anchor=east] at (1.4,\y) {\n};
  % пищалка
  \draw[thick,accent] (1.4,1) -- (4.2,1) to[buzzer,l=пищалка,color=accent] (6.4,1) -- (7.4,1) -- (7.4,-1.7);
  \node[font=\sffamily\scriptsize,anchor=south] at (4.2,0.55) {+};
  % кнопка
  \draw[thick,warn] (1.4,-0.5) -- (4.2,-0.5) to[normally open push button,l=кнопка,color=warn] (6.4,-0.5)
    -- (7.4,-0.5);
  \fill[black!70] (7.4,-0.5) circle (1.5pt);
  \draw[thick,black!70] (7.4,-1.7) -- (1.4,-1.7);
  \node[font=\sffamily\scriptsize,anchor=west,text=black!70] at (7.6,-1.1) {общий GND};
\end{circuitikz}
\caption{Подключение пассивной пьезопищалки к \pin{21} и кнопки к \pin{5}. Резисторы не нужны:
пищалка потребляет микроамперы, а подтяжку кнопки включает \code{INPUT\_PULLUP}}
\end{figure}

\subsection{Функция \texttt{tone()}}

\code{tone(вывод, частота, длительность)} запускает звук и сразу возвращается: программа продолжает
работать, а звук выключится сам. Без третьего аргумента звук длится, пока не вызвать \code{noTone(вывод)}.

\codecap{Звуковой сигнал при нажатии кнопки}
\codeinput{l24_beep}

\subsection{Мелодия без остановки программы}

Простейший способ сыграть мелодию --- цикл с \code{delay()} после каждой ноты. Но пока играет музыка,
такая программа глуха к кнопкам. Правильный способ уже знаком по уроку 21: мелодия --- это автомат,
который помнит номер текущей ноты и время её начала.

\codecap{Проигрыватель мелодии на \code{millis()}}
\codeinput{l24_melody}

Мелодия записана таблицей «частота --- длительность», нулевая частота означает паузу. Между нотами
программа делает тишину в \SI{30}{ms}: без неё две одинаковые ноты подряд сольются в одну длинную.

\begin{tip}
Громкость пьезопищалки заметно зависит от частоты: у большинства она максимальна в районе
\SIrange{2}{4}{kHz}. Для сигналов тревоги выбирайте частоты из этого диапазона, а для мелодий мириться
с тихим звуком помогает корпус-резонатор --- положите пищалку на пустую коробку.
\end{tip}

\subsection{Что дальше: настоящий звук}

Пищалка выдаёт только прямоугольный сигнал. Чтобы воспроизводить речь и музыку, у ESP32-S3 есть
интерфейс I\textsuperscript{2}S: по трём проводам он передаёт цифровой звук на усилитель (например,
MAX98357A) или принимает его с цифрового микрофона (INMP441). Это тема для итогового проекта.

\subsection{Самостоятельный опыт: музыкальная клавиатура}

Подключите матричную клавиатуру из урока 23. Клавиши «1»--«8» играют ноты от C4 до C5, пока клавиша
нажата; при отпускании звук прекращается.

\subsection{То же на ESPUI: ноты в списке на телефоне}

Вместо клавиатуры можно взять телефон: список нот и кнопка «Играть». Это хороший способ
проверить таблицу частот урока на слух и увидеть, как выбор из списка превращается в частоту.

Панели ESPUI мы уже делали в уроке 5; здесь два шага. Сеть называется \code{S3-Notes}
(пароль \code{lesson-39-demo}), адрес панели --- \code{http://192.168.4.1/}.

\subsubsection{Шаг 1. Список нот}

Сначала только выбор: список из восьми нот и подпись с частотой выбранной ноты. Пищалка пока не нужна:
проверяем, что панель открывается и список работает. Обработчик списка только записывает номер ноты
в \code{std::atomic}, а подпись меняет \code{loop()}.

\codecap{Шаг 1. Select ноты и подпись с частотой}
\codeinput[\footnotesize]{l24_espui_notes_s1}

\begin{itemize}[leftmargin=*]
  \item \code{ControlType::Select} --- выпадающий список, а каждая строка в нём --- элемент \code{Option};
  \item в \code{sender->value} приходит номер строки текстом, \code{toInt()} превращает его в число;
  \item \code{loop()} обновляет подпись только когда номер изменился, а не сто раз в секунду.
\end{itemize}

\subsubsection{Шаг 2. Кнопка «Играть» и звук}

Добавляем кнопку и пищалку между GPIO21 и GND. Нажатие только поднимает флаг \code{playRequest}, а
\code{loop()} забирает его, запускает \code{tone()} и запоминает момент, когда звук надо выключить.

\codecap{Шаг 2. Кнопка «Играть» включает звук на 400 мс}
\codeinput[\footnotesize]{l24_espui_notes}

\begin{itemize}[leftmargin=*]
  \item \code{playRequest.exchange(false)} --- прочитать флаг и сразу сбросить, чтобы звук включился один раз на щелчок;
  \item \code{tone(BUZZER\_PIN, f)} без длительности включает звук, а \code{stopAt} хранит время его конца;
  \item условие с \code{millis()} выключает звук позже, поэтому программа ни разу не останавливается на \code{delay()}.
\end{itemize}

\begin{trouble}
\begin{itemize}[leftmargin=*]
  \item \textbf{Вместо ноты слышны щелчки} --- пищалка активная: ей нужен постоянный уровень, а не частота.
  \item \textbf{Звук очень тихий} --- частота далека от резонанса пищалки; попробуйте \SIrange{2}{4}{kHz}.
  \item \textbf{Перестал работать ШИМ на другом выводе} --- \code{tone()} использует ту же периферию LEDC;
        задавайте звук и яркость на разных выводах и не вызывайте для вывода пищалки \code{ledcAttach()}.
\end{itemize}
\end{trouble}

\begin{summary}
\begin{itemize}[leftmargin=*]
  \item Высота звука --- это частота сигнала; удвоение частоты даёт октаву.
  \item Для мелодий нужна пассивная пищалка и функции \code{tone()} и \code{noTone()}.
  \item Мелодию играет автомат на \code{millis()}, а не цикл с \code{delay()}.
\end{itemize}
\end{summary}

\begin{selfcheck}
\begin{enumerate}
  \item Чем пассивная пищалка отличается от активной и как их различить?
  \item Во сколько раз отличаются частоты одной и той же ноты в соседних октавах?
  \item Почему мелодию не стоит играть циклом с \code{delay()}?
\end{enumerate}
\end{selfcheck}

\begin{tasks}
\begin{enumerate}
\item \lv{1} Сыграйте ноту ля (\SI{440}{Hz}) и измерьте её длительность по секундомеру.
\item \lv{1} Составьте таблицу частот для семи нот одной октавы и проиграйте гамму.
\item \lv{2} Запишите в таблицу свою мелодию из 10--15 нот и проиграйте её.
\item \lv{2} Проиграйте мелодию без остановки программы: кнопка должна реагировать во время звучания.
\item \lv{3} Сделайте сигнализацию: при нажатии кнопки звучит сирена --- частота плавно растёт от 800 до \SI{1800}{Hz} и падает обратно.
\item \lv{4} Сделайте метроном: потенциометр задаёт темп от 40 до 200 ударов в минуту, каждый четвёртый удар выше остальных.
\end{enumerate}
\end{tasks}
